\subsection{Very ampleness excludes the K3 alternative}

Very ampleness gives a stronger conclusion than ampleness in the
smallest equality case. A hyperplane section would have genus too
large for its degree unless the surface lay in projective
four-space. The normal bundle then gives a numerical contradiction.
We first give the curve estimate in arbitrary characteristic.

\begin{lemma}\label{lem:small-degree-curve-bound}
Let $B\subset\PP^n$, $n\geq4$, be a smooth integral nondegenerate
curve over an algebraically closed field. If its degree is
$d\leq9$ and its genus is at least two, then
\[
 g(B)\leq\max(d-4,0)+\max(d-7,0)<d+1.
\]
\end{lemma}

\begin{proof}
We use Kadets's geometric monodromy theorem: the sectional
monodromy of a nondegenerate integral nonstrange curve in
$\PP^n$ is $n$-transitive
\cite[Proposition~2.8 and proof, pp.~4--5]{Kadets}.
A curve is strange if all its tangent lines pass through one
point. The theorem holds in arbitrary characteristic.

The curve $B$ is nonstrange. Indeed, suppose all tangent lines
passed through $q\in\PP^n$. The twisted Euler sequence is
\[
 0\longrightarrow\OO_B(-1)\longrightarrow\OO_B^{n+1}
 \longrightarrow T_{\PP^n}|_B(-1)\longrightarrow0.
\]
A vector representing $q$ gives a section of the last bundle
which is nonzero generically. Its image in
$N_{B/\PP^n}(-1)$ is zero: away from $q$ the section is the
direction of the joining line, which is tangent to $B$.
The tangent--normal sequence consists of vector bundles, so it
factors globally through $T_B(-1)$, also when $q$ lies on $B$.
This would be a nonzero section of a line bundle of degree
$2-2g(B)-d<0$, a contradiction.

Let $\Gamma$ be a geometric generic hyperplane section, a set
of $d$ distinct points in $\PP^{n-1}$. Distinctness follows
also directly: hyperplanes through a fixed tangent line have
dimension $n-2$, so the incidence of tangent hyperplanes has
dimension at most $n-1$ in the $n$-dimensional dual space.
The points of $\Gamma$ span their hyperplane. To see this, use
\[
 0\longrightarrow\mathcal I_B\longrightarrow\mathcal I_B(1)
 \longrightarrow\mathcal I_{\Gamma/\PP^{n-1}}(1)
 \longrightarrow0.
\]
The hyperplane equation is a nonzerodivisor on $B$, which gives
exactness. Nondegeneracy gives $H^0(\mathcal I_B(1))=0$.
Moreover $H^1(\mathcal I_B)=0$, from
$H^0(\OO_{\PP^n})\simeq H^0(\OO_B)$ and
$H^1(\OO_{\PP^n})=0$. Hence
$H^0(\mathcal I_{\Gamma/\PP^{n-1}}(1))=0$.

In particular $d\geq n$ and some $n$ points of $\Gamma$ are
linearly independent. Linear dependence is invariant under the
sectional monodromy group. Its $n$-transitivity therefore makes
every $n$ distinct points independent; thus $\Gamma$ is in
linear general position. This uses no assertion of uniform
position in higher degree.

Write $h_\Gamma(m)$ for the homogeneous-coordinate Hilbert
function. Linear general position implies
\[
 h_\Gamma(m)\geq\min\{d,m(n-1)+1\}\qquad(m\geq1).
\]
Here is a direct interpolation argument. Choose any
$t=\min\{d,m(n-1)+1\}$ of the points and distinguish one.
Partition the other $t-1$ points into $m$ groups of at most
$n-1$ points. Each group lies on a hyperplane avoiding the
distinguished point, by linear general position. The product
of these hyperplane equations vanishes at the other selected
points and is nonzero at the distinguished point. Repeating
this for each selected point proves that evaluation has rank
$t$.

Let $R$ be the homogeneous coordinate ring of $B$, and let
$\ell$ define the generic hyperplane. Multiplication by
$\ell$ is injective. The quotient $R/\ell R$ surjects onto
the homogeneous coordinate ring of $\Gamma$, so
\[
 h_B(m)-h_B(m-1)\geq h_\Gamma(m).
\]
Since $h_B(0)=1$ and $h_B(m)=dm+1-g(B)$ for sufficiently
large $m$, summation gives
\[
 g(B)\leq\sum_{m\geq1}(d-h_\Gamma(m))
 \leq\sum_{m\geq1}\max\{d-m(n-1)-1,0\}.
\]
The possible saturation kernel in $R/\ell R$ has the required
sign; projective normality is not assumed. For $n\geq4$ and
$d\leq9$, the terms with $m\geq3$ vanish. The two remaining
terms are bounded by $\max(d-4,0)+\max(d-7,0)$.
For $d=4,\ldots,9$ these values are respectively
$0,1,2,3,5,7$, each smaller than $d+1$; degrees below four
were already excluded by $d\geq n$.
\end{proof}

\begin{proposition}\label{prop:very-ample-numerical-obstruction}
There is no smooth projective integral surface $X$ over an
algebraically closed field, of any characteristic, with a very
ample divisor $H$ having a smooth integral member, such that
\[
 K_X\equiv H,\qquad \chi(\OO_X)=1,
 \qquad 1\leq H^2\leq9.
\]
Here $\equiv$ denotes numerical equivalence.
\end{proposition}

\begin{proof}
Embed $X$ by the complete system $|H|$ in $\PP^r$ and put
$d=H^2$. A smooth integral member $C\in|H|$ is a hyperplane
section with
\[
 2g(C)-2=(K_X+H)\cdot H=2d,
 \qquad g(C)=d+1.
\]
It is nondegenerate in its hyperplane: the kernel of
$H^0(X,H)\to H^0(C,H|_C)$ is the one-dimensional space
generated by its defining section. If $r\geq5$,
Lemma~\ref{lem:small-degree-curve-bound} contradicts this
genus. Thus $r\leq4$.

An embedded surface requires $r\geq2$. For $r=2$, the image
is all of $\PP^2$, where $K_X=-3H$, a contradiction. For
$r=3$, the image is a degree-$d$ hypersurface and adjunction
gives $K_X=(d-4)H$. Numerical equivalence forces $d=5$.
The sequence
\[
 0\longrightarrow\OO_{\PP^3}(-5)
 \longrightarrow\OO_{\PP^3}\longrightarrow\OO_X
 \longrightarrow0
\]
then gives $h^0(\OO_X)=1$, $h^1(\OO_X)=0$ and
$h^2(\OO_X)=4$, by the cohomology of projective space
\cite[Tag 01XT]{Stacks}. Thus $\chi(\OO_X)=5$, again a
contradiction.

It remains to exclude $r=4$. Write $i:X\hookrightarrow\PP^4$
and $E=N_{X/\PP^4}$. This is a regular immersion of smooth
schemes, with tangent--normal sequence
\[
 0\longrightarrow T_X\longrightarrow T_{\PP^4}|_X
 \longrightarrow E\longrightarrow0.
\]
The Euler sequence gives $c_1(T_{\PP^4}|_X)=5H$ and
$c_2(T_{\PP^4}|_X)=10H^2$. As $c_1(T_X)=-K_X$, the first
two Chern terms give
\[
 c_1(E)=5H+K_X,\qquad
 \deg c_2(E)=10d+5K_X\cdot H+K_X^2-c_2(X).
\]
Noether's formula $K_X^2+c_2(X)=12\chi(\OO_X)$ follows
in every characteristic from Grothendieck--Riemann--Roch for
the structure morphism \cite[Tag 02UO]{Stacks}; it is an
equality of numerical Chern numbers. Consequently
\[
 c_2(X)=12-d,\qquad \deg c_2(E)=17d-12.
\]

Self-intersection gives the other value $\deg c_2(E)=d^2$.
For precision, base-change $i$ along itself. The fibre product
is $X$, its map to $X$ is the identity, and its excess
conormal bundle is $E^\vee$. The excess-intersection formula
\cite[Tag 0FV9]{Stacks} gives
\[
 i^!i_*[X]=c_2(E)\cap[X].
\]
Let $h=c_1(\OO_{\PP^4}(1))$. Since
$i_*[X]=d h^2\cap[\PP^4]$, commutation of $i^!$ with
Chern classes and $i^![\PP^4]=[X]$ give
\[
 i^!i_*[X]=d H^2\cap[X].
\]
These identities are Lemmas~42.59.7--8 of
\cite[Tag 0FEZ]{Stacks}; smoothness supplies the dimension
and Cohen--Macaulay hypotheses. Taking degrees gives $d^2$.
Thus $d^2-17d+12=0$, whose discriminant $241$ lies strictly
between $15^2$ and $16^2$. It has no integer root.
\end{proof}

\subsection{The distinguished cover in the almost-ruled case}

The remaining almost-ruled alternative determines the particular
cover defined by the torsion line bundle, rather than only some
purely inseparable cover of the surface.

\begin{proposition}\label{prop:distinguished-ruled-cover}
Let $X$ be a smooth projective integral surface over an
algebraically closed field $k$ of characteristic $p>0$.
Let $Y\subset X$ be a smooth integral ample divisor, and let
$M\ne\OO_X$ have order $p$ with $M|_Y\simeq\OO_Y$.
Let $N$ be the normalization of the $\mu_p$-torsor defined by
a trivialization of $M^{\otimes p}$.
If a ruled surface $R$ admits a dominant purely inseparable
rational map of degree $p$ to $X$, then
\[
 k(R)=k(N)\quad\text{as extensions of }k(X).
\]
In particular $N$ is birationally ruled, and every smooth
projective model of its function field has Kodaira dimension
$-\infty$. A resolution of the rational map from a smooth
projective ruled model gives a proper birational morphism
to $N$ over $X$. The surface $X$ itself is not birationally ruled.
Moreover, if $k(N)=A(t)$ is a ruling presentation, then
\[
 k(X)\cap A=A^p,\qquad k(X)A=k(N).
\]
\end{proposition}

\begin{proof}
Choose local frames $e_i$ of $M$ and write $u_i\in\OO_X^*$
for their $p$th powers under the chosen trivialization.
On an overlap, $u_i=g_{ij}^p u_j$, so the forms $du_i/u_i$
glue to a regular global form $\omega$ on $X$.
It is nonzero. Indeed, for a rational frame $e$, let
$u\in K^*$ be its $p$th power, where $K=k(X)$.
If $du=0$, the perfect-field differential criterion
\cite[Tag~031W]{Stacks} gives $u=v^p$ with $v\in K$.
The rational section $e/v$ has $p$th power one.
Its divisor has every coefficient zero, so normality makes
it a nowhere-vanishing global section of $M$, a contradiction.
Thus $u\notin K^p$, and the torsor has the degree-$p$ field
\[
 E=K(u^{1/p}).
\]
Its locally free algebra injects into that generic field,
so the torsor is integral. Its normalization equals the
normalization of $X$ in $E$, since integrality is transitive.
This normalization is finite: apply the Nagata finiteness
criterion to $\Spec E\to X$
\cite[Lemma~29.54.14, Tag~0BAK]{Stacks}.

The image of $\omega|_Y$ in $\Omega^1_{Y/k}$ is zero.
A trivialization of $M|_Y$ has $p$th power differing from
the chosen trivialization by a global unit on the proper
integral curve $Y$, hence by an element of $k^*$.
Its logarithmic differential is zero.
The conormal sequence is
\[
 0\longrightarrow\OO_Y(-Y)\longrightarrow
 \Omega^1_{X/k}|_Y\longrightarrow\Omega^1_{Y/k}
 \longrightarrow0.
\]
Since $Y^2>0$, the left-hand line bundle has no nonzero
global section. Hence $\omega|_Y$ vanishes in the full
restricted cotangent bundle. There is consequently a
nonzero section
\[
 \theta\in H^0(X,\Omega^1_{X/k}(-Y)),
 \qquad s_Y\theta=\omega,
\]
where $s_Y$ is the canonical section of $\OO_X(Y)$.

Write the ruled function field as $k(C)(t)$, with $C$ a
smooth projective integral curve, and replace $R$ by
$R_0=C\times\PP^1$. Resolve $R_0\dashrightarrow X$ by
point blowups, obtaining $S\to R_0$ and $g:S\to X$
\cite[Tag~0C5H]{Stacks}. Point blowups of a regular surface
remain regular \cite[Lemma~54.3.2, Tag~0AGP]{Stacks}; since
$k$ is perfect, $S$ is smooth. It is projective, and
$[k(S):K]=p$.

Only finitely many ruling fibres contain images of the
blowup centres. A general fibre $F$ of $S\to C$ is therefore
an unmodified projective line, with
\[
 \Omega^1_{S/k}|_F\simeq\OO_{\PP^1}\oplus\OO_{\PP^1}(-2).
\]
The restriction $g|_F$ is nonconstant: a general ruling
fibre meets the dense quasi-finite locus of $g$ and cannot
be contracted. Ampleness of $Y$ gives
$e_F=\deg(g^*\OO_X(Y)|_F)>0$, without any separability
assumption on this curve map.

Pullback by the morphism $g$ gives a regular section
\[
 \eta\in H^0(S,\Omega^1_{S/k}(-g^*Y)).
\]
It restricts on each general ruling fibre to a section of
\[
 \OO_{\PP^1}(-e_F)\oplus\OO_{\PP^1}(-e_F-2),
\]
and is therefore zero there. These fibres fill a dense open
of the integral surface $S$, so $\eta=0$ and $g^*\omega=0$.
This uses positivity only on the general fibre; $g^*Y$
need not be ample on $S$.

At the generic point, $du/u=0$ in $\Omega^1_{k(S)/k}$.
The same perfect-field criterion gives a $p$th root of $u$
in $k(S)$, and hence a $K$-embedding $E\hookrightarrow k(S)$.
Both extensions have degree $p$, so they are equal.
The $p$th root is unique, and changing the frame or the
trivialization changes $u$ only by a $p$th power; thus this
is the distinguished extension defined by $M$.

The resulting birational map $S\dashrightarrow N$ over $X$
extends to a morphism. For an affine open $U=\Spec B$ of
$X$, let $B'$ be the integral closure of $B$ in $E$.
Every $b'\in B'$ is integral over every local ring
$\OO_{S,s}$ for $s\in g^{-1}(U)$, and belongs to it by
normality. These regular functions give
$g^{-1}(U)\to\Spec B'$ and glue over $X$.
The resulting map is birational by the field equality and
proper by its closed graph and properness of $g$.
Thus $N$ has a ruled function field.

For any smooth projective model $S'$ of that field, resolve
the birational map $R_0\dashrightarrow S'$ in the same way.
The resulting ruled blowup has no positive pluricanonical
sections: its canonical bundle has degree $-2$ on every
general unmodified ruling fibre. Pullback by the birational
morphism injects the pluricanonical sections of $S'$ into
those of the blowup. Hence every positive plurigenus of
$S'$ is zero.

If $X$ itself were birationally ruled, the argument with
$\theta$ could instead be applied to a birational map
$R_0\dashrightarrow X$. It would give $g^*\omega=0$, although
the map of function fields and hence of their differentials
is an isomorphism. This contradicts $\omega\ne0$.

Finally write $L=k(N)=A(t)$. The purely inseparable
degree-$p$ extension $L/K$ has exponent one, so
$L^p\subset K$ and $A^p\subset A\cap K\subset A$.
Since $A$ is the function field of a curve over the perfect
field $k$, one has $[A:A^p]=p$. For example, a separating
generator $x$ gives $[A:k(x)]=[A^p:k(x^p)]$, and the tower
through $k(x^p)\subset k(x)$ proves the assertion.
If $A\subset K$, then
$A(t^p)\subset K\subset A(t)$ and the degree-$p$ equality
forces $K=A(t^p)$, contradicting non-ruledness of $X$.
Thus $A\cap K=A^p$. The strict intermediate compositum
$KA$ in the prime-degree extension $L/K$ is all of $L$.
\end{proof}

The field intersection does not assert that the rational ruling
map on $X$ extends everywhere, or that its fibres are smooth.
The normalization may be singular away from the inverse image
of $Y$; resolving those singularities changes an ample pullback
into a divisor having degree zero on exceptional curves.


\begin{corollary}\label{cor:very-ample-equality-almost-ruled}
Let $X$ be a smooth projective integral surface over an
algebraically closed field of characteristic $p>0$, and let
$Y\subset X$ be a smooth integral very ample divisor. If a
nontrivial line bundle $M$ of order $p$ satisfies
\[
 M|_Y\simeq\OO_Y,\qquad 2g(Y)-2=pY^2,
\]
then $X$ is almost ruled: a ruled surface admits a dominant
purely inseparable rational map of degree $p$ to $X$.
The normalization of the $\mu_p$-torsor defined by $M$ is
birationally ruled, with the same extension of $k(X)$ as
every such degree-$p$ ruled cover. Its smooth projective
models have Kodaira dimension $-\infty$, whereas $X$ is
not birationally ruled.
\end{corollary}

\begin{proof}
The nonnegative terms in the identity of
Proposition~\ref{prop:normalized-mup-conductor} force $a=1$
and both defects to vanish. If $X$ were not almost ruled,
Proposition~\ref{prop:normalized-mup-ekedahl} would give
$p=2$, $\chi(\OO_X)=1$, $\omega_X^2\simeq\OO_X(2Y)$
and $1\leq Y^2\leq9$. In particular $K_X\equiv Y$.
Apply Proposition~\ref{prop:very-ample-numerical-obstruction}
with $H=Y$ and its smooth member $Y$ to obtain a contradiction.
The conclusions about the distinguished cover and non-ruledness
of $X$ now follow from
Proposition~\ref{prop:distinguished-ruled-cover}.
\end{proof}

This excludes the K3 alternative when the divisor is very ample.
The distinguished normal cover in the almost-ruled case is therefore
birationally ruled. The following subsections analyze its positive-genus
and rational rulings, and Theorem~\ref{thm:strict-picard-all} excludes
the smallest very-ample equality in every positive characteristic. For a merely
ample divisor, the K3 existence question stated above remains.
